Upon successful completion of this course, students will be able to:
\begin{ul}
    \item Analyze opposing views on government intervention in the economy, as supported by different macroeconomic schools of thought.
    \item Identify and assess possible sources of bias in macroeconomic measures.
    \item Consider the pros and cons of using macroeconomic statistics like GDP or unemployment to characterize quality of life.
    \item Use supply and demand models to analyze cause-effect patterns graphically and numerically.
    \item Calculate key macroeconomic statistics using raw data.
    \item Apply relevant economic theories to predict the long- and short-run effects of changes to monetary and fiscal policy.
    \item Analyze whether the pursuit of self-interest can promote the social interest.
    \item Discuss the social and ethical considerations involved in the pursuit of efficiency versus equality.
    \item Characterize the ways in which today's economic choices can help or hurt future generations.
\end{ul}

This course also fulfills the following university-wide institutional competencies:
\begin{ul}
    \item \textbf{Analytical Thinking:} The ability to reason, interpret, analyze, and solve problems from a wide array of authentic contexts.
    \item \textbf{Critical Thinking:} The ability to pursue and comprehensively evaluate information before accepting or establishing a conclusion, decision, or action.
    \item \textbf{Social Awareness \& Responsibility:} The capacity to understand the interdependence of people, communities, and self in a global society.
\end{ul}
